% !TeX root = ../../Book5.tex
% 纲要标题：### 2.2 正交关系
% 纲要位置：工作记录/计划纲要.md，第 4088 行。
% 纲要细目（写作范围）：
% * 矩阵系数正交关系
% * 特征标内积与行正交
% * 不可约特征标构成类函数空间的基与列正交


\section{正交关系}
\label{b5:sec:0202}

正交关系来自对线性映射取平均。平均后得到交织算子，Schur 引理迫使它在不同不可约表示之间为零，在相同表示上为标量。选取矩阵单位就把这个结论变成矩阵条目及特征标的求和公式。


\subsection{矩阵系数正交关系}
\label{b5:subsec:0202:01}

选定表示 \(\rho\) 的基后，每个矩阵条目给出函数 \(g\mapsto\rho(g)_{ia}\)，称为 \(\rho\) 的矩阵系数。它属于全函数空间 \(\mathbb C^G\)，以下使用\eqref{b5:eq:character-inner-product}在该空间上定义的内积。

\begin{theorem}[矩阵系数正交]\label{b5:thm:matrix-coefficient-orthogonality}
	设 \(G\) 为有限群，\(\rho\) 与 \(\sigma\) 为 \(G\) 的不可约有限维复表示，维数分别为 \(d_\rho,d_\sigma\)。在各自的不变 Hermitian 内积下选定正交归一基，将两表示写成酉矩阵。若它们不同构，则对 \(1\le i,a\le d_\sigma\)、\(1\le j,b\le d_\rho\) 有
	\[
		\frac1{|G|}\sum_g\sigma(g)_{ia}\overline{\rho(g)_{jb}}=0.
	\]
	若 \(\rho=\sigma\) 是同一个 \(d\) 维不可约表示，并对两者采用同一组正交归一基，则对 \(1\le i,a,j,b\le d\) 有
	\[
		\frac1{|G|}\sum_g\rho(g)_{ia}\overline{\rho(g)_{jb}} =\frac1d\delta_{ij}\delta_{ab}.
	\]
\end{theorem}
\begin{proof}
	对 \(T:V_\rho\to V_\sigma\) 定义
	\[
		P(T)=\frac1{|G|}\sum_g\sigma(g)T\rho(g^{-1}).
	\]
	重排求和表明 \(P(T)\) 是从 \(V_\rho\) 到 \(V_\sigma\) 的交织算子；若 \(T\) 本来就是交织算子，则平均式中每一项均为 \(T\)，所以 \(P\) 投影到交织算子空间。两表示不同构时，由 Schur 引理得 \(P(T)=0\)。当 \(\rho=\sigma\) 且采用同一组基时，Schur 引理说明这个投影将全部线性算子送到标量直线 \(\mathbb C I\)。写 \(P(T)=cI\)，共轭不改变迹，故 \(cd=\operatorname{tr}T\)，因而
	\[
		P(T)=\frac{\operatorname{tr}T}{d}I.
	\]

	取矩阵单位 \(T=E_{ab}:V_\rho\to V_\sigma\)，即把 \(V_\rho\) 的第 \(b\) 个基向量送到 \(V_\sigma\) 的第 \(a\) 个基向量，其余基向量送到零。在同一表示上，非对角矩阵单位的迹为零，故平均后消失；每个对角矩阵单位的迹为一，故都变成 \(I/d\)。指标 \(a,b\) 指定被平均的矩阵单位，而 \(i,j\) 指定平均所得矩阵中要读取的条目。平均式的这个条目为
	\[
		P(E_{ab})_{ij} =\frac1{|G|}\sum_g\sigma(g)_{ia}\rho(g^{-1})_{bj} =\frac1{|G|}\sum_g\sigma(g)_{ia}\overline{\rho(g)_{jb}},
	\]
	最后一步使用 \(\rho\) 的酉性。不同构时该条目为零；同一表示、同一组基时，\(\operatorname{tr}E_{ab}=\delta_{ab}\)，该条目为 \(\delta_{ab}\delta_{ij}/d\)，得到两式。
\end{proof}

这还说明同一不可约表示的 \(d^2\) 个矩阵系数线性无关：它们互相正交，且每个范数平方为 \(1/d>0\)。矩阵系数一般不是类函数，只有取对角元之和才得到特征标。

\subsection{特征标内积与行正交}
\label{b5:subsec:0202:02}

\begin{theorem}\label{b5:thm:character-hom-inner-product}
	设 \(G\) 为有限群，\(V,W\) 为有限维复表示。则
	\[
		\langle\chi_V,\chi_W\rangle_G =\dim_{\mathbb C}\operatorname{Hom}_G(V,W).
	\]
	特别地，不可约特征标两两正交且范数为一。
\end{theorem}
\begin{proof}
	记 \(H=\operatorname{Hom}_{\mathbb C}(V,W)\)。在这个向量空间上，对共轭作用 \(T\mapsto\rho_W(g)T\rho_V(g^{-1})\) 取平均。这里每个线性映射 \(T:V\to W\) 都是 \(H\) 的一个向量，平均算子 \(P\) 则是 \(H\) 上的线性算子。它为投影，像恰为不动空间 \(\operatorname{Hom}_G(V,W)\)，故
	\[
		\operatorname{tr}_{H}P=\dim_{\mathbb C}\operatorname{Hom}_G(V,W).
	\]
	这个迹是在 \(H\) 上取的；即使 \(V=W\)，取迹的对象仍是 \(H\) 上的 \(P\)，而非 \(V\) 上的 \(P(T)\)。例如，若 \(S\) 是二维不可约表示，共轭平均在四维空间 \(\operatorname{End}_{\mathbb C}(S)\) 上投影到 \(\mathbb C I\)，所以 \(\operatorname{tr}_{\operatorname{End}_{\mathbb C}(S)}P=1\)；但 \(P(I)=I\)，在 \(S\) 上取迹却得 \(\operatorname{tr}_{S}P(I)=2\)。

	由\cref{b5:prop:hom-tensor-invariants}，\(H\) 的表示与 \(W\otimes_{\mathbb C}V^*\) 同构。张量迹公式因此给出共轭作用在 \(H\) 上的迹为
	\[
		\chi_W(g)\chi_V(g^{-1})=\overline{\chi_V(g)}\chi_W(g).
	\]
	对 \(g\) 求平均，利用迹的线性性就得到 \(\operatorname{tr}_H P=\langle\chi_V,\chi_W\rangle_G\)，即所需等式。不可约时，再用\cref{b5:lem:schur}和\cref{b5:cor:schur-algebraically-closed}，同构时 Hom 维数为一，不同构时为零。
\end{proof}

\begin{corollary}\label{b5:cor:character-irreducibility}
	设 \(G\) 为有限群，\(V\) 为 \(G\) 的有限维复表示，\(S_i\) 为两两不同构的不可约复表示、\(m_i\) 为非负整数。若 \(V\) 的不可约分解为 \(\bigoplus_i S_i^{\oplus m_i}\)，则
	\[
		\langle\chi_{S_i},\chi_V\rangle_G=m_i,\qquad \langle\chi_V,\chi_V\rangle_G=\sum_i m_i^2.
	\]
	因此非零 \(V\) 不可约当且仅当后一数为一；两个复表示同构当且仅当特征标相同。
\end{corollary}
\begin{proof}
	Maschke 定理保证存在分解；迹在直和上相加，所以 \(\chi_V=\sum_i m_i\chi_{S_i}\)。逐项用刚证的正交关系即可得到公式。重数唯一确定分解的同构类型，而 \(\sum_i m_i^2=1\) 恰好表示只有一个重数为一，其他为零。
\end{proof}

这些是特征标表的行正交关系。列正交还需要知道不可约特征标已经填满类函数空间，不能仅从一组正交向量直接假定它们是基。

\subsection{不可约特征标与类函数空间的基}
\label{b5:subsec:0202:03}

\begin{theorem}\label{b5:thm:irreducible-character-basis}
	设 \(G\) 为有限群。其全部不可约复特征标组成复类函数空间 \(\operatorname{CF}(G)\) 关于特征标内积 \(\langle\ ,\ \rangle_G\) 的标准正交基。因而不可约复表示的同构类型数等于共轭类数。
\end{theorem}
\begin{proof}
	由\cref{b5:lem:irreducibles-in-regular}，不可约类型只有有限个；正交性已经给出线性无关。若类函数 \(f\) 与全部不可约特征标正交，要把这些内积等式变成群代数元素在所有不可约表示上的零作用。为此用 \(f\) 的值作系数，并安排逆元，使作用的迹成为特征标内积，令
	\[
		z_f=\sum_{g\in G}f(g^{-1})g\in\mathbb C[G].
	\]
	对群代数元素 \(z=\sum_g a_gg\) 与 \(h\in G\)，重新编号给出
	\[
		hzh^{-1}=\sum_g a_g hgh^{-1}=\sum_g a_{h^{-1}gh}g.
	\]
	比较群元素基的系数，\(hz=zh\) 等价于对所有 \(g\) 有 \(a_{h^{-1}gh}=a_g\)。因此 \(z\) 与全部群元素交换，恰好意味着系数在每个共轭类上常值。因此 \(z_f\) 位于中心；让 \(f\) 遍历全部类函数，上述构造给出复向量空间同构
	\[
		\Phi:\operatorname{CF}(G)\longrightarrow Z(\mathbb C[G]),\qquad f\longmapsto z_f.
	\]
	其逆把 \(\sum_g a_gg\) 送到函数 \(h\mapsto a_{h^{-1}}\)。若 \(\mathbf1_C\) 是共轭类 \(C\) 的示性函数，则 \(\Phi(\mathbf1_C)=\sum_{g\in C^{-1}}g\)；取逆将共轭类相互排列，所以各共轭类的示性函数组成的基，被 \(\Phi\) 送到由元素 \(\sum_{g\in C}g\) 组成的中心基，其中 \(C\) 遍历全部共轭类。这些求和元素将在\secref{b5:sec:0203}中记为 \(z_C\)。

	现在利用中心性：\(z_f\) 的作用与群作用交换，故在任一不可约 \(S_i\) 上由 Schur 引理作用为 \(a_iI\)；这样迹为零便能推出整个算子为零。具体地，换元 \(h=g^{-1}\) 及 \(\chi_i(h^{-1})=\overline{\chi_i(h)}\) 给出
	\[
		(\dim S_i)a_i=\sum_g f(g^{-1})\chi_i(g) =|G|\langle\chi_i,f\rangle_G=0.
	\]
	故在每个不可约上作用为零。正则表示是这些不可约的直和，所以 \(z_f\) 在正则表示上也为零；作用于单位基向量得到 \(z_f=0\)，于是所有系数 \(f(g^{-1})=0\)。所以正交补为零，给定正交组确为整个类函数空间的基。
\end{proof}

\begin{theorem}[列正交]\label{b5:thm:row-column-orthogonality}
	设 \(G\) 为有限群，全部不可约复特征标为 \(\chi_1,\ldots,\chi_r\)。对 \(g,h\in G\)，有
	\[
		\sum_{i=1}^r\chi_i(g)\overline{\chi_i(h)} =
		\begin{cases}
			|C_G(g)|,&g,h\;\text{共轭},\\
			0,&g,h\;\text{不共轭}.
		\end{cases}
	\]
\end{theorem}
\begin{proof}
	上一基定理保证不可约特征标数等于共轭类数。取各共轭类 \(C_j\) 的代表 \(g_j\)，令类大小 \(c_j=|G|/|C_G(g_j)|\)。类示性函数也是类函数空间的一组基，且 \(\langle\mathbf1_{C_j},\mathbf1_{C_k}\rangle_G=(c_j/|G|)\delta_{jk}\)，所以将它们归一化为
	\[
		b_j=\sqrt{\frac{|G|}{c_j}}\,\mathbf1_{C_j},\qquad 1\le j\le r,
	\]
	就得到另一组标准正交基。按内积在第二变量线性的约定，\(\chi_i\) 在这组基下的展开为 \(\chi_i=\sum_j\langle b_j,\chi_i\rangle_G b_j\)。将每个 \(\chi_i\) 的坐标排成一行，得到方阵
	\[
		U_{ij}=\langle b_j,\chi_i\rangle_G =\sqrt{\frac{c_j}{|G|}}\,\chi_i(g_j).
	\]
	归一化因子因而来自两组标准正交基的坐标转换。行正交给出 \((U\overline U^{\mathsf T})_{ik}=\langle\chi_k,\chi_i\rangle_G=\delta_{ik}\)，即 \(U\overline U^{\mathsf T}=I\)。方阵的右逆也为左逆，故 \(\overline U^{\mathsf T}U=I\)。写出其第 \(j,k\) 个条目，得到
	\[
		\delta_{jk} =\frac{\sqrt{c_jc_k}}{|G|} \sum_{i=1}^r\overline{\chi_i(g_j)}\chi_i(g_k).
	\]
	当 \(j=k\) 时，这给出
	\[
		\sum_{i=1}^r|\chi_i(g_j)|^2 =\frac{|G|}{c_j}=|C_G(g_j)|.
	\]
	当 \(j\ne k\) 时，前面的实系数为正，故求和为零；对该零等式取复共轭，得到
	\[
		\sum_{i=1}^r\chi_i(g_j)\overline{\chi_i(g_k)}=0.
	\]
	特征标取值与中心化子阶都在共轭类上不变，因此这两式推广到任意 \(g,h\) 就是所述结论。
\end{proof}

行正交计算表示重数，列正交计算中心化子阶。单位元的一列特别给出
\[
	\sum_i\chi_i(1)^2=|G|.
\]
这一平方和也可以从正则表示的具体分解得到，下一节将保留那个更有结构的信息。
